\documentclass[11pt,a4paper]{article}
\usepackage[utf8]{inputenc}\usepackage[T1]{fontenc}\usepackage[italian,english]{babel}
\usepackage{amsmath,amssymb,amsthm}\usepackage[margin=2.2cm]{geometry}\usepackage{longtable,booktabs,array,calc}
\usepackage{listings,xcolor}\usepackage{hyperref}\usepackage{url}
\providecommand{\tightlist}{\setlength{\itemsep}{0pt}\setlength{\parskip}{0pt}}
\lstset{basicstyle=\ttfamily\scriptsize,breaklines=true,frame=single,captionpos=b,columns=fullflexible,keepspaces=true,inputencoding=utf8,extendedchars=true}
\title{Proof Pursuit --- Problem 1, part 5\\ \large prove, decisioni degli agenti, codice verificato, letteratura}
\author{Team Proof Pursuit (Researcher: T. Tumini; Referee: gabundos; Creative: M. Renzi; gio3r3gio)}
\date{\today}
\begin{document}\maketitle\tableofcontents
\section*{Come leggere questo rapporto}
Per ogni cella: la richiesta ufficiale, lo stato dichiarato dal team (mai ``risolto'' senza revisione umana), la consegna
con la prova, poi la traccia delle decisioni degli agenti (Researcher: famiglia e motivo; Referee: verdetto ed errore
fatale; Creative: quando e perché è intervenuto) e il codice su cui il tentativo poggia con l'esito della riesecuzione
fidata. DIMOSTRATO, CITATO, VERIFICATO SU INTERVALLO FINITO e CONGETTURA restano distinti. L'architettura del sistema
è descritta nel README del repository.
\section{Problem 1 — Angles between lines (Fejes Toth)}
\subsection{Cella 5 (8 punti) — stato: parziale}
\subparagraph{\texorpdfstring{Parte 5 --- \(d+2\) lines in
\(\mathbb{R}^d\)}{Parte 5 --- d+2 lines in \textbackslash mathbb\{R\}\^{}d}}\label{parte-5-d2-lines-in-mathbbrd}

\textbf{Punteggio:} 8 points · \textbf{Valutazione:} Judged

The case \(N = d+2\) in every dimension, with the same conjectured
optimum: the \(d\) coordinate axes, two of them repeated. Prove that for
every \(d \ge 2\), any \(d+2\) lines in \(\mathbb{R}^d\) satisfy \[
S \;\le\; \left( \binom{d+2}{2} - 2 \right) \frac{\pi}{2}.
\]

\textbf{Enunciato protetto dato agli agenti (state.json).} \texttt{Prove that for every d \textgreater{}= 2, any d+2 lines through the origin in R\textasciicircum{}d satisfy S \textless{}= (C(d+2,2) - 2) pi/2.}

\textbf{Evidenza registrata in STATUS.md.} triage; bozza parziale in inglese in submission/ (parziali.py)

\subsubsection{Consegna (prova)}
\paragraph{Problema 1 --- Parte 5 --- bozza di consegna
PARZIALE}\label{problema-1-parte-5-bozza-di-consegna-parziale}

\textbf{Stato dichiarato: PARZIALE.} Nessuna soluzione completa: qui
sotto ciò che è stabilito, la formalizzazione, la posizione rispetto
alla letteratura e ciò che resta aperto. Nulla è dichiarato dimostrato
oltre quanto scritto.

\subparagraph{1. Risultato e ambito}\label{risultato-e-ambito}

\textbf{Richiesta ufficiale.} \#\# Parte 5 --- \(d+2\) lines in
\(\mathbb{R}^d\)

\textbf{Punteggio:} 8 points · \textbf{Valutazione:} Judged

The case \(N = d+2\) in every dimension, with the same conjectured
optimum: the \(d\) coordinate axes, two of them repeated. Prove that for
every \(d \ge 2\), any \(d+2\) lines in \(\mathbb{R}^d\) satisfy \[
S \;\le\; \left( \binom{d+2}{2} - 2 \right) \frac{\pi}{2}.
\]

\textbf{Cosa consegniamo.} Nessuna prova. Riformulazione in termini di
deficit (\(\sum\delta_{ij}\ge\pi\)) e sanity check dei valori.

\subparagraph{2. Dimostrazione}\label{dimostrazione}

La tesi equivale a: \(d+2\) versori in \(\mathbb R^d\) hanno deficit
totale \(\ge\pi\). Se la parte 3 si ottiene per riduzione a una catena
(deficit \(\ge\pi/2\)), la via naturale è iterare la riduzione su due
circuiti indipendenti della dipendenza lineare (rango \(\le d\) fra
\(d+2\) vettori: lo spazio delle relazioni ha dimensione \(\ge2\)). Non
sviluppato.

\subparagraph{3. Verifica: istruzioni, dipendenze,
tempi}\label{verifica-istruzioni-dipendenze-tempi}

Codice disponibile (Python 3, libreria standard; ogni script gira in
meno di un minuto): - \texttt{problema-1/esperimenti/lines\_Rd.py} -
\texttt{problema-1/esperimenti/p1\_somma\_angoli.py}

\subparagraph{4. Fonti e contributo}\label{fonti-e-contributo}

Letteratura arXiv (ricerca deterministica
\texttt{tools/cerca\_letteratura.sh}, abstract letti, non usata come
prova): - arXiv:1801.07837v1 --- On the Fejes Tóth Problem about the Sum
of Angles Between Lines (Dmitriy Bilyk, Ryan W Matzke, 2018); solo
abstract letto. - arXiv:2007.08698v2 --- On Fejes Tóth's conjectured
maximizer for the sum of angles between lines (Tongseok Lim, Robert J.
McCann, 2020); solo abstract letto. - arXiv:1204.3850v1 --- Simple
Agents Learn to Find Their Way: An Introduction on Mapping Polygons
(Jérémie Chalopin, Shantanu Das, Yann Disser, 2012); solo abstract
letto. Non risulta noto (Bilyk--Matzke {[}1801.07837{]}; Lim--McCann
{[}2007.08698{]}).

\subparagraph{5. Limiti e parti
irrisolte}\label{limiti-e-parti-irrisolte}

Dipende dalla parte 3.

\subsubsection{Decisioni degli agenti}
\paragraph{Run \texttt{p1\_c5}}

\subsubsection{Codice su cui poggia il tentativo}
Nessun codice nei tentativi.

\section{arXiv literature consulted}
\begin{thebibliography}{99}
\bibitem{1801.07837v1} Dmitriy Bilyk, Ryan W Matzke. \emph{On the Fejes Tóth Problem about the Sum of Angles Between Lines}. arXiv:1801.07837v1 (2018). \url{http://arxiv.org/abs/1801.07837v1}
\bibitem{2007.08698v2} Tongseok Lim, Robert J. McCann. \emph{On Fejes Tóth's conjectured maximizer for the sum of angles between lines}. arXiv:2007.08698v2 (2020). \url{http://arxiv.org/abs/2007.08698v2}
\bibitem{1204.3850v1} Jérémie Chalopin, Shantanu Das, Yann Disser, Matúš Mihalák, Peter Widmayer. \emph{Simple Agents Learn to Find Their Way: An Introduction on Mapping Polygons}. arXiv:1204.3850v1 (2012). \url{http://arxiv.org/abs/1204.3850v1}
\end{thebibliography}
\end{document}
